\subsection{7. Derselbe Grundzustand besitzt eine kontrollierte geladene
Antwort}\label{derselbe-grundzustand-besitzt-eine-kontrollierte-geladene-antwort}

Für die präzisen Energie- und Gewichtsschranken verwenden wir den
engeren Bereich

\[\adjustbox{max width=\linewidth}{$\displaystyle 0<|g|/\Delta\le1/10000.$}\]

\subsubsection{7.1 Der gesamte Rest wird als konvergente Resolvente
kontrolliert}\label{der-gesamte-rest-wird-als-konvergente-resolvente-kontrolliert}

P projiziere auf alle bosonfreien Zustände, \(\Pi=I-P\),
\(D=\Delta N_b|_\Pi\), \(X=\sum(b^\dagger P_A+P_A^\dagger b)\), und
\(C=\Pi X P\). Dann \(C^\dagger C=Q_{\rm ges}|_P\). Für
\(-Mg^2/\Delta\le E\le0\) schreibe

\[\adjustbox{max width=\linewidth}{$\displaystyle \Pi(H-E)\Pi=D^{1/2}(I+Z)D^{1/2},\qquad
Z=gD^{-1/2}\Pi X\Pi D^{-1/2}-ED^{-1}.$}\]

Die Bosonzahlschichten geben

\[\adjustbox{max width=\linewidth}{$\displaystyle \|Z\|\le2\sqrt M\,|g|/\Delta+M g^2/\Delta^2<1/100.$}\]

Der Einbosonblock des linearen X-Terms ist null. Deshalb ist die
Fehlergrenze des Feshbach-Operators gegenüber
\(-(g^2/\Delta)Q_{\rm ges}\) nicht nur erster Ordnung in \(g/\Delta\),
sondern

\[\adjustbox{max width=\linewidth}{$\displaystyle \left\|g^2C^\dagger[\Pi(H-E)\Pi]^{-1}C-\frac{g^2}{\Delta}Q_{\rm ges}\right\|
\le\frac{g^2}{\Delta}M\left[\frac{M}{10^8}+\frac{(1/100)^2}{1-1/100}\right]
=\frac{g^2}{\Delta}\frac{119008}{515625}<\frac{g^2}{4\Delta}.$}\]

Hier wurde der vollständige Neumannrest summiert. Der
Schur-/Feshbach-Schritt ist elementare Operatorblockrechnung; ein
allgemeiner Literaturrahmen ist
\href{https://arxiv.org/abs/2105.02058}{Dusson, Sigal und Stamm}. Die
konkreten Konstanten und Quellenoperatoren oben wurden hier abgeleitet.

Insbesondere

\[\adjustbox{max width=\linewidth}{$\displaystyle -1920\frac{g^2}{\Delta}\le E_0<-1919.75\frac{g^2}{\Delta},\qquad
\operatorname{gap}(H)>1.5\frac{g^2}{\Delta}.$}\]

Im N=256-Sektor ist der Schurraum eindimensional. Der exakte
Grundzustandswert ist daher die eindeutige negative Lösung einer
skalaren Gleichung,

\[\adjustbox{max width=\linewidth}{$\displaystyle \boxed{E_0=-g^2\langle F|C^\dagger[\Pi(H-E_0)\Pi]^{-1}C|F\rangle.}$}\]

Das ist ein Fixpunkt des \textbf{tatsächlichen} gewählten
Hamiltonoperators. Es ist nicht die Umbenennung eines Compilerpolynoms
in eine physische Bewegungsgleichung. Die Resolvente und die
rekonstruierte Zustandskomponente besitzen konvergente Reihen; die
genaue Dezimalenergie wird hier nicht als berechnet ausgegeben.

\subsubsection{7.2 Zwei isolierte
Entnahmelinien}\label{zwei-isolierte-entnahmelinien}

Im bosonfreien N=255-Sektor gibt es genau ein Loch. Normalordnen
beziehungsweise die lokale Casimir-Identität liefern dort exakt

\[\adjustbox{max width=\linewidth}{$\displaystyle Q_{\rm ges}|_{\text{ein Loch}}=MI-15(U^\dagger U)\otimes I_{64}.$}\]

Der geometrische Fermionclock hat vier verschiedene Eigenphasen
\(\exp(i(2j+1)\pi/4)\). Pro Phase trägt der Einlochraum genau eine
irreduzible 64-dimensionale innere Darstellung. Weil H beide Symmetrien
erhält, ist sein energieabhängiger Schuroperator in jedem dieser vier
Blöcke skalar. Er hat jeweils eine negative Nullstelle: Seine Ableitung
nach E ist strikt negativ, und der Einloch-Kopplungsoperator ist
injektiv. Die Realität von H und des Clockoperators paart konjugierte
Phasen. Damit ergeben sich \textbf{genau zwei verschiedene niedrige
Energielinien}, jeweils mit Vielfachheit 128. Die getrennten Intervalle
unten schließen eine Zusammenlegung der beiden Linien aus.

Setze

\[\adjustbox{max width=\linewidth}{$\displaystyle a_-=15(1-1/\sqrt2),\qquad a_+=15(1+1/\sqrt2).$}\]

Aus derselben uniformen Schurrestschranke und der
Grundenergieeinschließung folgen die Entnahmeenergien

\[\adjustbox{max width=\linewidth}{$\displaystyle \boxed{(a_\pm-1/2)\frac{g^2}{\Delta}
<\epsilon_\pm<(a_\pm+1/4)\frac{g^2}{\Delta}.}$}\]

In Einheiten \(g^2/\Delta\) ergeben sich die nach außen gerundeten
Intervalle \((3.8933,4.6435)\) und \((25.1066,25.8567)\). Dies sind
analytische Einschließungen, keine Polzentralwerte. Alle übrigen
N=255-Zustände liegen im absoluten Spektrum mindestens bei
\(\Delta/2-2Mg^2/\Delta>0\); ihr Abstand zum negativen Grundwert ist
entsprechend größer. Der Additionssektor N=257 erfüllt ebenfalls diese
positive Untergrenze.

\subsubsection{7.3 Ein ursprüngliches, clockaufgelöstes Fermion sieht
die
Linie}\label{ein-urspruxfcngliches-clockaufgeluxf6stes-fermion-sieht-die-linie}

Für den Grundzustand und die vier niedrigen Einlochblöcke ergibt die
Resolventenschranke dieselbe untere Überlappung mit dem jeweiligen
bosonfreien Referenzzustand:

\[\adjustbox{max width=\linewidth}{$\displaystyle Z\ge Z_*=\frac{408375}{408383}>0.99998.$}\]

Ein normierter Clock-Eigenmodus \(a_{\theta,r}\) ist eine unitäre
Linearkombination der \textbf{ursprünglichen} f-Moden, kein neu
eingeführtes Kompositfeld. Da \(a_{\theta,r}\) die Bosonzahl nicht
ändert und Operatornorm eins hat, ist sein Übergang auf die zugehörige
niedrige Linie mindestens

\[\adjustbox{max width=\linewidth}{$\displaystyle \boxed{|\langle\Psi_{\theta,r},a_{\theta,r}\Omega\rangle|^2
\ge(2Z_*-1)^2=\frac{166763606689}{166776674689}>0.99992.}$}\]

Beweis: Der bosonfreie Anteil der Amplitude ist mindestens
\(\sqrt{Z_0Z_\theta}\), der Betrag der übrigen Kreuzung höchstens
\(\sqrt{(1-Z_0)(1-Z_\theta)}\). Ihre Differenz ist mindestens
\(2Z_*-1\). Da \(\|a_{\theta,r}\Omega\|^2\le1\), gilt dieselbe
Untergrenze auch für den Anteil am normierten Entnahmespektrum.

Somit trägt die zugehörige Linie \textbf{über 99.992 Prozent} für einen
clockaufgelösten Modus. Ein ursprünglicher einzelner Orts-/Kopienmodus
enthält hingegen beide Clock-Paare und sieht im Allgemeinen beide
niedrigen Linien. Für ihn wird keine einzelne Linie mit diesem Gewicht
behauptet.

Hier gehören Quelle, Grundzustand und geladene Antwort wirklich zu
demselben H. Die native Herstellung eines Fermion-Eingriffs oder eines
Messgeräts ist dadurch weiterhin nicht konstruiert.

\subsection{8. Was wir bei der Einfachheit möglicherweise falsch
verlangt
haben}\label{was-wir-bei-der-einfachheit-muxf6glicherweise-falsch-verlangt-haben}

\subsubsection{8.1 Quadratische Erzeugung ist nicht vollständige
dynamische
Steuerbarkeit}\label{quadratische-erzeugung-ist-nicht-vollstuxe4ndige-dynamische-steuerbarkeit}

Der alte Prüfer \codeword{v111\_quadratic\_transport.py} zeigt korrekt,
dass quadratische Cliffordwörter auf dem 16-dimensionalen geraden
Fockraum eine volle \textbf{assoziative} Matrixalgebra aufspannen.
Daraus folgt nicht, dass zeitabhängige quadratische Hamiltonoperatoren
jede unitäre Operation erzeugen. Ihre dynamische Lie-Algebra hat
Dimension 45, nicht 255.

Die Ergänzung muss dennoch keine riesige Maschine sein: Gewährt man alle
quadratischen Generatoren und \textbf{einen} nichtverschwindenden reinen
quartischen Majoranaterm, erzeugen die Kommutatoren sämtliche 210
quartischen Richtungen. Zusammen mit den 45 quadratischen Richtungen
entsteht \(\mathfrak{su}(16)\). Der Prüfer kontrolliert die 210
erreichbaren Vierermengen und die exakte Hilbert-Schmidt-Unabhängigkeit
aller 255 Richtungen.

Das ist ein möglicher minimaler Operationsbaustein, keine Behauptung,
dieser steuerbare Term liege bereits nativ vor. Die fünf inneren
Cliffordslots sind außerdem nicht mit den 64 Teilchenmoden der W-Bank
gleichzusetzen. Die gekoppelte W-Quelle erzeugt zwar quartische
effektive Terme; deren Gleichsetzung mit einem frei steuerbaren
Quartikgenerator des inneren 16er-Raums ist gerade noch nicht gezeigt.
Die Unterscheidung zwischen assoziativer Erzeugung und dynamischer
Lie-Algebra entspricht dem etablierten Kontrollrahmen, etwa
\href{https://doi.org/10.1038/s41534-024-00900-2}{Wiersema et al.,
Klassifikation dynamischer Lie-Algebren}.

\subsubsection{8.2 Ein eindeutiges mikroskopisches Netz ist nicht
zwingend
nötig}\label{ein-eindeutiges-mikroskopisches-netz-ist-nicht-zwingend-nuxf6tig}

Die Petersen-/Prisma-Gegenmodelle beweisen eine endliche
Unterbestimmtheit, aber keine Verschiedenheit aller denkbaren
großskaligen Grenztheorien. Unterschiedliche Mikroregeln können
derselben Universalitätsklasse angehören. Das ist ein wesentlicher
Grund, nicht immer feinere Graph-Eindeutigkeit zu verlangen, bevor eine
physische Grenzrechnung begonnen wird. Siehe den ursprünglichen
RG-Rahmen von
\href{https://theory.tifr.res.in/~tridib/ReferenceMaterial/WilsonKogut.pdf}{Wilson
und Kogut}.

Diese Entlastung ist keine automatische Rettung: Auf der einfachen
unendlichen Zyklus-/Kettenfortsetzung hat die N=2-Antwort

\[\adjustbox{max width=\linewidth}{$\displaystyle B(k)=8+4\cos k,\qquad
E_-(k)=\frac{\Delta-\sqrt{\Delta^2+4g^2B(k)}}2,$}\]

also \(E_-(k)-E_-(0)=2g^2k^2/\sqrt{\Delta^2+48g^2}+O(k^4)\) und keinen
linearen Dirackegel an dieser Bandkante. Dies ist nur die
N=2-Bandrechnung, \textbf{nicht} die Anregung des vollständigen großen
Grundzustands und kein universeller Ausschluss des Modells.

Auch ein antiperiodischer Zyklus mit immer mehr gerader Knotenzahl hat
\(s_{\min}=1-\cos(\pi/L)\to0\). Der neue endliche Grundzustandsbeweis
ist deshalb kein bereits gleichmäßig gappender Thermodynamikbeweis. Für
den Standardmodell-/Raumzeitanspruch müssen Zustand, Anregung,
Skalierung und Lokalität gemeinsam geprüft werden.

\subsection{9. Was die alte E8-Kaskade wirklich
beisteuert}\label{was-die-alte-e8-kaskade-wirklich-beisteuert}

Die gewünschte Archivdatei ist ein LaTeX-Umschlag einer ausdrücklich als
bestmögliche PDF-Textextraktion bezeichneten 81-seitigen Fassung. Einige
Formeln und Figuren sind in dieser Darstellung beschädigt oder fehlen.
Aus solchen Schäden werden hier keine mathematischen Gegenbeweise
konstruiert.

Die lesbare Skalenregel ist

\[\adjustbox{max width=\linewidth}{$\displaystyle D_n=60-2n,\qquad
\varphi_n=\varphi_0e^{-\gamma_0}(D_n/58)^\lambda,\qquad
\lambda=\frac{\gamma_0}{\log(248/60)},\quad\gamma_0=0.834.$}\]

Sie liefert eine logarithmisch organisierte Skalenleiter, nicht allein
eine Hamiltonentwicklung. In Verhältnissen fällt der gemeinsame
Vorfaktor heraus, aber die Abhängigkeit von \(\gamma_0\) bleibt über
\(\lambda\) erhalten:

\[\adjustbox{max width=\linewidth}{$\displaystyle \frac{\partial}{\partial\gamma_0}\log\frac{\varphi_m}{\varphi_n}
=\frac{\log(D_m/D_n)}{\log(248/60)}\ne0\quad(m\ne n).$}\]

Die aktuellen Kaskadenkapitel zitieren einen
Krümmungs-/Glättungsselektor. Der dort angegebene Prüfer
\codeword{v5\_e8\_cascade.py} kontrolliert Endpunkte, Dimensionen und
Summen, aber \textbf{nicht} diese Auswahl von \(\gamma_0\) oder
\(\lambda\). Eine lexikografisch eindeutige Kette unter erklärten
Kriterien ist zudem nicht schon eine physische Erklärung dieser
Kriterien.

Noch klarer ist die Grenze in der alten Callan-Symanzik-Passage: Die
dort ausgeschriebene, bei den behaltenen Ordnungen positive Betafunktion

\[\adjustbox{max width=\linewidth}{$\displaystyle \beta(\alpha)=\frac{b_1}{2\pi}\alpha^2+A c_3^2\alpha^3+\cdots$}\]

besitzt bei \(b_1,A,c_3>0\) in genau dieser Trunkierung keine positive
nichttriviale Nullstelle. Das anschließend genannte Compilerkubikum kann
daher nicht einfach als dieser RG-Fixpunkt ausgegeben werden. Das
widerlegt nicht den algebraischen Wert des Compilerkubikums; es sperrt
die Gleichsetzung zweier verschiedener Bedeutungen von Fixpunkt.

Die Archivabschnitte zur Chiralität setzen bereits \(M_4\) mal eine
innere Faser, einen Spinor-/Randprojektorvertrag und einen Fluss an. Sie
enthalten also keine bislang übersehene vollständige Erzeugung von
3+1D-Raumzeit aus W. Der produktive alte Hinweis ist stattdessen:
\textbf{Umlauf, Vorzeichen und Skalierung gehören zusammen}, dürfen aber
nicht als derselbe Zeitbegriff behandelt werden.

\subsection{10. Einstein, Noether und Dirac als
Arbeitsprinzipien}\label{einstein-noether-und-dirac-als-arbeitsprinzipien}

\textbf{Noether:} Die erhaltene Ladung \(N_f+2N_b\) ist hier kein
dekoratives Etikett. Sie reduziert den globalen Grundzustandsvergleich
auf untere, mittlere und obere Ladungssektoren und macht den kurzen
Beweis möglich.

\textbf{Dirac:} Die Paarantwort ist nicht das vollständige fermionische
Objekt. Die Hebung, ihr Vorzeichen und die Gradierung müssen vor der
Quadratisierung bewahrt werden. Das ist eine konkrete
Double-Cover-Aussage, nicht die Behauptung, eine beliebige Quadratwurzel
liefere bereits den relativistischen Diracoperator.

\textbf{Einstein:} Der nächste physische Test betrifft gemeinsame
Kovarianz, Kausalstruktur und Energie desselben Modells. Passende
Konstanten und ein formal kovarianter Zielraum ersetzen diese gemeinsame
Realisierung nicht.

\textbf{Die einfache Leitidee:} Nicht noch mehr Zahlen sammeln, sondern
die kleine \textbf{gradierte Quelle vor dem Informationsverlust}
festhalten; daraus Zustand und Antworten ableiten; erst dann prüfen,
welche räumliche Grenztheorie sie tatsächlich trägt.

\subsection{11. T1-T8 bleiben einzeln
sichtbar}\label{t1-t8-bleiben-einzeln-sichtbar}

{\def\LTcaptype{none} % do not increment counter
\begin{longtable}[]{@{}
  >{\raggedright\arraybackslash}p{(\linewidth - 4\tabcolsep) * \real{0.3333}}
  >{\raggedright\arraybackslash}p{(\linewidth - 4\tabcolsep) * \real{0.3333}}
  >{\raggedright\arraybackslash}p{(\linewidth - 4\tabcolsep) * \real{0.3333}}@{}}
\toprule\noalign{}
\begin{minipage}[b]{\linewidth}\raggedright
Gate
\end{minipage} & \begin{minipage}[b]{\linewidth}\raggedright
Was diese Runde tatsächlich beiträgt
\end{minipage} & \begin{minipage}[b]{\linewidth}\raggedright
Was weiter fehlt
\end{minipage} \\
\midrule\noalign{}
\endhead
\bottomrule\noalign{}
\endlastfoot
T1 Auswahl & Zwei Spin-Lifts, ein Clock-Intertwiner, ein bedingter
Reflexionsselektor und Zustandsselektion in einer kleinen Klasse &
Ursprüngliche gemeinsame Vertex-/Kantenreflexion, Kopienregel und
Kopplungsvertrag \\
T2 markierte Seam & Die verlorene fermionische Hebungsinformation wird
präzisiert & Renormiertes Half-Charge-Feld,
Energie-/Adjungiertenkontrolle und tatsächliche
E8-Seam-Identifikation \\
T3 gemeinsamer 3+1D-Ursprung & Ein konkret zusammengesetzter endlicher
Parent & Raumdimension, lokale Algebren, Skalierung und derselbe
physische Parent für alle Sektoren \\
T4 chirales Maß & Kein neuer vollständiger Feldadapter & Chirales
Eich-/Weylmaß, Anomalien, Spiegelentkopplung \\
T5 Kontinuum & Vollständiger endlicher Fock-Grundsatz mit
Bosonrestkontrolle & Gleichmäßiger wechselwirkender Grenzübergang und
relativistische Anregungen \\
T6 Parameter & Zwei kontrollierte Modelllinien aus W und Clock &
Herleitung von g/Delta und der Verbindung zu beobachteten Massen,
Kopplungen und Neutrinos \\
T7 Gravitation & Kein neuer Spin-2-Nachweis & Quantisierter masseloser
Spin 2 und universelle Kopplung aus demselben Parent \\
T8 Zustand/Instrumente & Eindeutiger globaler Modellgrundzustand und
geladene Antwort auf ihm & Physische Herkunft der Quelle sowie native
Präparation, autonome Kontrolle und beschreibbares Record \\
\end{longtable}
}

\subsection{12. Die nächsten entscheidenden
Prüfungen}\label{die-nuxe4chsten-entscheidenden-pruxfcfungen}

\begin{enumerate}
\def\labelenumi{\arabic{enumi}.}
\tightlist
\item
  \textbf{Ursprüngliche Quelle statt weiterer freier Graphen:} Eine
  konkrete Abbildung von TFPTs Seam-/Cliffordoperationen auf das
  gewählte q-Feld angeben, die CAR, Clock, W-Vertex und Zustand
  gemeinsam erhält. Abschnitt 15 präzisiert den Angriffspunkt: die
  gemeinsame Vertex-/Kantenreflexion. Übereinstimmung der Clockordnung
  oder einer abstrakten dihedralen Relation allein reicht nicht.
  Scheitert die Zustandsverträglichkeit, ist das ein tatsächlicher
  Ausschluss dieses Adapters.
\item
  \textbf{Den bewiesenen Modellbereich erweitern:} Für denselben
  antiperiodischen Kandidaten den Kopplungsweg bis \(g/\Delta=1/20\)
  kontrollieren. Dabei Ladungswechsel, zusätzliche niedrige Zustände und
  Polgewichte verfolgen. Der vorhandene Satz darf nicht über seinen
  kleinen Bereich hinaus zitiert werden.
\item
  \textbf{Physische Universalitätsklasse statt vollständiger
  UV-Eindeutigkeit:} Eine vorab festgelegte Verfeinerungsfolge nehmen
  und dieselbe geladene Antwort auf ihrem jeweiligen Grundzustand
  untersuchen. Entscheidend sind kleine Impulse, Skalierung, lokale
  Antwort und ein gemeinsamer relativistischer Feldtyp. Erst ein
  erfolgreicher solcher Test trägt die Arbeit an chiraler Materie und
  Gravitation.
\end{enumerate}

Keine dieser drei Prüfungen wird hier als bereits bestanden bezeichnet.
Der Gewinn ist, dass die erste Zustands-/Antwortkette für einen
konkreten Lift nun ausgeführt ist und ihr Beweis nicht von einer
astronomischen Matrix abhängt.

\subsection{13. Reproduktion und
Lesereihenfolge}\label{reproduktion-und-lesereihenfolge}

\codeword{verify\_spin\_lift.py} prüft native Tensorpins, die
vollständige benötigte Casimir-Grundlage, die N3-Reduktion,
Vorzeichenklassen, unabhängige CAR-Kontrollen, exakte Momente, rationale
Restkonstanten sowie die Operations- und Archivpräzisierungen.
\codeword{replay.py} wiederholt diese Rechnung normal und optimiert und
führt den historischen v480-Zeugen separat aus. Der vollständige
analytische Grundzustands-/Antwortbeweis steht in den Abschnitten 5-7;
bloßes Zählen grüner Bedingungen ersetzt ihn nicht.

Das Prüfprotokoll benennt ausdrücklich, welche Ergebnisse endliche
Gleichheiten, welche analytische Implikationen und welche noch
Herkunftshypothesen sind. Alte Dokumente und Originaleingaben werden
nicht überschrieben. Die aktualisierte Hauptfassung enthält diese neue
Ebene und die vollständige vorherige Hauptfassung v1.6.9 als
gekennzeichneten historischen Bestand; das kurze Update ersetzt sie
nicht.

